\section{Introduction}
\IEEEPARstart{C}{ooperative} target protection requires defenders to redistribute their effort as vehicle and target populations change. An approaching attacker may already be covered by a teammate, while another threatens an exposed asset. Actions therefore depend on relationships among attackers, targets, and defenders, interpreted from the perspective of the defender choosing an action. A policy trained in small engagements must handle different distributions of these relationships when deployed with more vehicles or targets. Reusing one policy across these populations is a concrete requirement for automated cooperative protection.

Entity-based policies provide a common interface to changing scenes. Set encoders and attention accommodate variable entity populations~\cite{zaheer2017deepsets,lee2019settransformer}, and REFIL encourages reuse of interaction patterns through randomized entity-wise value factorization~\cite{iqbal2021refil}. Supporting a larger input set does not itself establish effective coordination. The controller must develop entity context and extract information relevant to the receiving defender. Repeated relational attention offers one way to process that context~\cite{zambaldi2019relational,dehghani2019universal}.

A second question is which processing stages should supply the decision feature. Retaining intermediate representations permits aggregation across depth~\cite{xu2018jk}. Image-classification networks can change correct intermediate predictions into incorrect final predictions~\cite{kaya2019overthinking}; pure-attention rank-loss results also show that depth effects depend on architectural conditions, including residual and feedforward paths~\cite{dong2021attention}. These precedents motivate exposing intermediate contexts without assuming that the last context is always preferable. They do not establish those phenomena in a cooperative policy.

A third requirement concerns time. Public physical states need not reveal an attacker's assigned target, motivating memory across observations. Processing the current entity set and retaining temporal information are different operations. We place the relational reader before the recurrent utility network, allowing its query to depend on the current self representation while the downstream GRU retains history across physical decisions.

\textbf{Looped Entity Attention Fusion (LEAF)} combines these three roles. Each defender initializes a receiver-relative entity workspace, applies a shared contextual-update block, and reads its intermediate sets with one self-conditioned query. Learned weights fuse the reads, and a residual projection augments the REFIL feature before one GRU update. The query's previous-GRU-state input is zero and explicit count injection is disabled; downstream recurrence is retained. Fusion selects information from the computed rounds without changing execution depth. Training samples one through four rounds, and standard evaluation uses four. Figure~\ref{fig:framework} shows this separation of within-decision processing and cross-decision memory.

We evaluate frozen policies in HAD target protection against a fixed reactive attacker. Across three training seeds, LEAF has higher group mean reward $R=-D$, where $D$ is accumulated target damage, than five learning baselines and Rule in ID, target-count, vehicle-count, and joint extrapolation. Evaluation reaches five times the largest training vehicle population and four times the largest training target count, without retraining. Saved-design comparisons examine entity-context updating, query-history placement, and host integration; the complete variant results retain the context in which LEAF was selected.

Frozen interventions on a registered three-configuration subset connect these roles to the saved policies. Bypassing contextual updates while retaining the relation projection and GRU lowers all nine seed--configuration reward means. The first two reads dominate fusion, while rare two-round action changes concentrate near small greedy-action margins. Freezing fusion weights has mixed effects in the two $K=2$ configurations and lowers reward in all three seeds at $(30,30,12)$. Resetting downstream GRU memory lowers eight of nine means. These deployment-dependence measurements are interpreted separately from trained-design comparisons. The contributions are:
\begin{enumerate}
\item A cooperative control architecture that separates shared contextual entity processing, fusion of intermediate reads, and downstream temporal recurrence.
\item A frozen-policy population-generalization study across 24 configurations, using three training seeds and common evaluation scenes to compare LEAF with five learning baselines and a rule-based controller.
\item A mechanism analysis combining paired saved-design comparisons, frozen computation interventions, and a local action-margin condition to distinguish contextual-feature dependence, stage selection, and temporal memory use.
\end{enumerate}
